% Chapter 6

\chapter{总结与展望}

\section{研究工作回顾}

本文以大语言模型安全对齐漏洞为研究对象，围绕为什么越狱攻击能够成立、哪类攻击最为有效、规则型防御能够达到何种效果三个核心问题展开，在理论分析、攻击实验与防御设计三个层面进行了系统研究。

理论层面，本文从 SFT、RLHF 与 DPO 等后训练方法的内在机制出发，将安全对齐的失效归结为三个相互关联的结构性问题：有用性与安全性之间难以彻底消解的目标冲突、基于经验分布的对齐训练对拒答模式的学习而非对危险语义的理解，以及模型输出对上下文结构的依赖性\cite{rafailov2023dpo,hong2024orpo}\cite{ethayarajh2024kto,christiano2017preferences}。这三点共同说明，安全对齐本质上是一种软约束，而不是稳健的硬边界。

攻击实验层面，本文构建了覆盖九类静态攻击维度、自动化红队扩展与多轮越狱的统一实验框架，基于 \texttt{deepseek-chat}、\texttt{qwen2.5:1.5b} 和 \texttt{qwen2.5:3b} 三个模型开展了系统评测\cite{yu2024gptfuzzer,chao2024pair}\cite{mehrotra2024tap,liu2023autodan}。结果表明，认知误导类攻击在三个模型上均以显著优势领先其他维度，跨模型平均成功率达 12.64\%；多轮越狱将 \texttt{qwen2.5:3b} 上的突破率从 14.58\% 进一步推高至 43.33\%；而大量不确定样本（最高超过 78\%）提示，模型的安全边界在相当范围内仍处于不稳定状态，而非已经稳定失守或稳定拒答。

防御设计层面，本文针对上述攻击实验的分析结果，提出了输入层、交互层与输出层三层协同的防御方案，并在 \texttt{qwen2.5:3b} 上开展了对照实验\cite{rebedea2023nemo,phute2024llmselfdefense}\cite{chen2024struq,yi2023bipia}。结果表明，经过输入阈值、风险传播与输出替换机制的修正后，三层防御已能在单轮与多轮场景中提供稳定收益：单轮静态场景下，三层全开将越狱成功率从 2.23\% 压缩至 1.44\%，并将明确拒绝率提升至 26.27\%；多轮场景下，交互层成为关键控制环节，可将认知误导类攻击的累计突破率从 43.33\% 压缩至 28.33\%，三层全开则进一步压缩至 23.33\%。但整体而言，当前规则型防御对语义包装类攻击的覆盖仍然有限，这并非分层架构本身失效，而是规则型检测在面对语义驱动攻击时的共性边界。

\section{主要创新点}

在机理分析方面，本文将越狱现象归结为目标冲突、表层对齐与上下文依赖三类结构性问题，而不止于描述特定模板是否有效。这一框架使不同攻击维度的成功原因可以在统一视角下解释，也为防御设计提供了比规则库更稳定的理论依据。

在实验体系方面，九类攻击维度共享同一数据加载、执行流程与评判口径，结果之间可以直接横向比较。评测指标在显式成功率之外专门引入了边界模糊率，将那些兼含拒答信号与风险内容的中间态回复作为独立分析维度保留，避免在汇总统计中丢失这一层信号。自动化红队变体与多轮越狱实验则在同一基线下验证了静态边界不稳定性与动态突破潜力之间的对应关系。

在防御设计方面，三层方案将输入识别、跨轮次诱导追踪与输出内容检测分配到不同层，各层有效性边界的差异化分布决定了防御深度，而非由单层能力上限决定。交互层对多轮诱导的专属覆盖是现有单阶段防御方法所不具备的，也是本文方案在多轮场景中取得额外收益的直接来源。

在框架扩展性方面，本文将九类文本攻击维度与对抗性图像注入、排版视觉越狱、跨模态语义包装三类多模态攻击建立了基于机理的结构性对应，并讨论了分层防御框架各层向多模态场景延伸的实现路径。这部分工作并非实验验证，而是说明本文的维度划分逻辑在攻击机理层面具备跨模态适用性。

\section{工作局限}

本文的工作存在一定局限。实验覆盖的模型参数规模整体偏小，防御实验仅在 \texttt{qwen2.5:3b} 上开展，结论向更大规模商业模型的可推广性有待后续验证；多轮越狱分析主要集中于认知误导维度，工具调用、智能体任务链等更复杂场景尚未涉及。规则型防御对语义包装类攻击的覆盖范围受限，是此类方法的共性技术边界，引入语义层面的检测能力是后续改进的自然方向。

\section{后续研究方向}

本文对认知误导类攻击突破能力的分析表明，语义包装可能绕过规则检测，因此有必要进一步考察在防御端引入语义理解能力是否能够有效弥补这一盲区。从技术可行性来看，轻量级意图分类模型已较为成熟，在输入层或输出层嵌入此类能力并不依赖基础模型能力的根本性突破。未来研究可将这一方向作为重点拓展内容，以第五章的三层框架为基础，逐步以语义判断补充当前的规则匹配，并通过对照实验评估两者在误拦截率与覆盖率之间的权衡；在评测上，也可结合安全过度拒答测试集细化误报分析\cite{rottger2024xstest,wallace2024instructionhierarchy}。

从更长远的角度看，越狱攻击的边界将随着模型能力的扩展而同步扩展。当大语言模型开始承担工具调用、外部系统交互和多代理协同等任务时，安全问题将不再只是模型会不会输出危险文字，而会扩展到任务执行权限、外部操作链的约束传播等层面\cite{hagendorff2026jailbreakagents,toyer2024tensortrust}。进一步考虑多模态生成与跨模态红队场景时，这一问题还会延伸到图像生成和复合交互系统\cite{chin2024icer}。现有的文本对齐方法和规则型防御在这类场景中能够保留多少有效性，目前并不清楚，这是未来研究中最需要填补的空白之一。

\section{本章小结}

本章从理论分析、攻击实验、防御设计与框架扩展四个层面对全文主要工作进行了系统回顾，提炼了核心结论：安全对齐本质上是依赖统计分布的软约束，认知误导是当前最具威胁的攻击路径，边界模糊带是动态攻击的主要入口，分层防御能够提供可测量的协同收益但规则型框架的能力上限清晰可见。本章还归纳了本文在分析视角、实验体系、防御思路与框架扩展性四个方面的主要贡献：前三者共同构成从机制理解到实验验证再到防御设计的完整研究链条，第四项则通过多模态攻击维度的结构映射论证了本文框架的跨模态通用性。本章同时明确指出了当前工作在模型规模覆盖和多模态实验验证上的局限性，并展望了两个核心后续方向：在三层防御框架中引入语义级意图识别能力，以及开展面向视觉—语言模型的多模态越狱攻击实验验证。
